\documentclass[a4paper,11pt]{article}

% ============================================================
% PACKAGES
% ============================================================

\usepackage[utf8]{inputenc}
\usepackage[T1]{fontenc}
\usepackage[french]{babel}

\usepackage{amsmath}
\usepackage{amssymb}
\usepackage{bm}
\usepackage{mathtools}

\usepackage{geometry}
\usepackage{physics}

\usepackage{booktabs}
\usepackage{siunitx}

\geometry{
    left=2.5cm,
    right=2.5cm,
    top=2.5cm,
    bottom=2.5cm
}

% ============================================================
% COMMANDES
% ============================================================

\newcommand{\m}{\bm{m}}
\newcommand{\er}{\bm{e}_r}
\newcommand{\ephi}{\bm{e}_\varphi}
\newcommand{\ez}{\bm{e}_z}

%\newcommand{\dd}{\mathrm{d}}
\newcommand{\E}{\mathcal{E}}

% ============================================================
% DOCUMENT
% ============================================================

\title{
Minimisation numérique de l'énergie micromagnétique\\
d'une texture axisymétrique avec DMI de volume
}

\author{}
\date{}

\begin{document}

\maketitle

\tableofcontents

\newpage

% ============================================================
\section{Problème physique}
% ============================================================

On considère un cylindre ferromagnétique de rayon $R$ et d'épaisseur $t$.
L'axe du cylindre est pris suivant $Oz$.

Dans le cas considéré ici,

\begin{equation}
R = \SI{50}{nm},
\qquad
t = \SI{10}{nm}.
\end{equation}

On suppose que l'aimantation ne dépend pas de $z$ et que la texture est
invariante par rotation autour de $Oz$.

L'aimantation réduite est définie par

\begin{equation}
\m = \frac{\bm M}{M_s},
\end{equation}

où $M_s$ est l'aimantation à saturation. Par définition,

\begin{equation}
\boxed{|\m|=1.}
\end{equation}

On néglige ici l'énergie magnétostatique et on considère uniquement :

\begin{enumerate}
    \item l'énergie d'échange ;
    \item le DMI de volume ;
    \item l'anisotropie magnétocristalline uniaxiale.
\end{enumerate}

% ============================================================
\section{Paramétrisation de l'aimantation}
% ============================================================

On introduit un angle polaire $\theta(r)$ mesuré par rapport à l'axe $Oz$
et un angle d'hélicité $\phi$.

La texture est écrite

\begin{equation}
\boxed{
\m(r,\varphi)
=
\sin\theta(r)
\left[
\cos\phi\,\er
+
\sin\phi\,\ephi
\right]
+
\cos\theta(r)\,\ez.
}
\end{equation}

On a donc

\begin{equation}
m_r = \sin\theta\cos\phi,
\end{equation}

\begin{equation}
m_\varphi = \sin\theta\sin\phi,
\end{equation}

et

\begin{equation}
m_z = \cos\theta.
\end{equation}

L'angle $\phi$ est supposé constant dans tout le disque.

Les cas particuliers sont :

\begin{align}
\phi &= 0,\pi
&&\text{texture de Néel},\\
\phi &= +\frac{\pi}{2},-\frac{\pi}{2}
&&\text{texture de Bloch}.
\end{align}

Dans le programme numérique, $\phi$ est toutefois laissé libre afin que la
minimisation détermine elle-même l'hélicité de l'état fondamental.

% ============================================================
\section{Énergie micromagnétique}
% ============================================================

L'énergie totale s'écrit

\begin{equation}
E =
E_{\mathrm{ex}}
+
E_{\mathrm{DMI}}
+
E_{\mathrm{an}}.
\end{equation}

% ------------------------------------------------------------
\subsection{Énergie d'échange}
% ------------------------------------------------------------

L'énergie d'échange est

\begin{equation}
E_{\mathrm{ex}}
=
A
\int_V
(\nabla\m)^2
\,\dd V,
\end{equation}

où $A$ est la constante d'échange, en \si{J.m^{-1}}.

Pour la texture axisymétrique considérée,

\begin{equation}
(\nabla\m)^2
=
\left(
\frac{\dd\theta}{\dd r}
\right)^2
+
\frac{\sin^2\theta}{r^2}.
\end{equation}

Par conséquent,

\begin{equation}
E_{\mathrm{ex}}
=
2\pi t A
\int_0^R
\left[
r\left(\frac{\dd\theta}{\dd r}\right)^2
+
\frac{\sin^2\theta}{r}
\right]
\dd r.
\end{equation}

% ------------------------------------------------------------
\subsection{DMI de volume}
% ------------------------------------------------------------

On considère un DMI de volume de densité

\begin{equation}
w_{\mathrm{DMI}}
=
D\,\m\cdot(\nabla\times\m),
\end{equation}

où $D$ est la constante de DMI volumique, en \si{J.m^{-2}}.

Pour la texture considérée,

\begin{equation}
\boxed{
\m\cdot(\nabla\times\m)
=
\sin\phi
\left[
\frac{\dd\theta}{\dd r}
+
\frac{\sin\theta\cos\theta}{r}
\right].
}
\end{equation}

On obtient donc

\begin{equation}
E_{\mathrm{DMI}}
=
2\pi tD\sin\phi
\int_0^R
\left[
r\frac{\dd\theta}{\dd r}
+
\sin\theta\cos\theta
\right]
\dd r.
\end{equation}

% ------------------------------------------------------------
\subsection{Anisotropie}
% ------------------------------------------------------------

L'axe facile est supposé être $Oz$.

La densité d'énergie d'anisotropie est

\begin{equation}
w_{\mathrm{an}}
=
-K_u m_z^2,
\end{equation}

avec $K_u>0$.

Puisque

\begin{equation}
m_z=\cos\theta,
\end{equation}

on obtient

\begin{equation}
E_{\mathrm{an}}
=
-2\pi tK_u
\int_0^R
r\cos^2\theta
\,\dd r.
\end{equation}

% ============================================================
\section{Fonctionnelle totale}
% ============================================================

En regroupant les trois contributions,

\begin{equation}
\boxed{
E[\theta,\phi]
=
2\pi t
\int_0^R
\left\{
A
\left[
r\theta_r^2
+
\frac{\sin^2\theta}{r}
\right]
+
D\sin\phi
\left[
r\theta_r
+
\sin\theta\cos\theta
\right]
-
K_u r\cos^2\theta
\right\}
\dd r,
}
\end{equation}

où

\begin{equation}
\theta_r
=
\frac{\dd\theta}{\dd r}.
\end{equation}

Cette fonctionnelle est celle que l'on cherche à minimiser numériquement.

% ============================================================
\section{Adimensionnement}
% ============================================================

Pour améliorer le conditionnement numérique, on introduit la coordonnée
radiale réduite

\begin{equation}
\boxed{
x=\frac{r}{R}.
}
\end{equation}

Ainsi,

\begin{equation}
0\leq x\leq1,
\qquad
r=Rx,
\qquad
\dd r=R\dd x.
\end{equation}

De plus,

\begin{equation}
\frac{\dd\theta}{\dd r}
=
\frac{1}{R}
\frac{\dd\theta}{\dd x}.
\end{equation}

On définit les deux paramètres sans dimension

\begin{equation}
\boxed{
d=\frac{DR}{A}
}
\end{equation}

et

\begin{equation}
\boxed{
k=\frac{K_uR^2}{A}.
}
\end{equation}

L'énergie peut alors être écrite sous la forme

\begin{equation}
\boxed{
E
=
2\pi t A R\,\mathcal{E},
}
\end{equation}

où l'énergie réduite est

\begin{equation}
\boxed{
\mathcal{E}[\theta,\phi]
=
\int_0^1
\left[
x\theta_x^2
+
\frac{\sin^2\theta}{x}
+
d\sin\phi
\left(
x\theta_x
+
\sin\theta\cos\theta
\right)
-
kx\cos^2\theta
\right]
\dd x.
}
\end{equation}

Ici,

\begin{equation}
\theta_x
=
\frac{\dd\theta}{\dd x}.
\end{equation}

Le facteur global $2\pi tAR$ n'a aucune influence sur la position du
minimum. Le programme minimise donc directement $\mathcal{E}$.

% ============================================================
\section{Conditions aux limites}
% ============================================================

Au centre du cylindre, on impose

\begin{equation}
\boxed{
\theta(0)=0.
}
\end{equation}

L'aimantation au centre est donc

\begin{equation}
\m(0)=\ez.
\end{equation}

La valeur de $\theta(1)$ n'est pas imposée.

Elle doit être déterminée par la minimisation de l'énergie.

% ============================================================
\section{Condition de Brown généralisée}
% ============================================================

La variation de l'énergie fait apparaître un terme de bord.

La densité lagrangienne réduite est

\begin{equation}
\mathcal{L}
=
x\theta_x^2
+
\frac{\sin^2\theta}{x}
+
d\sin\phi
\left(
x\theta_x
+
\sin\theta\cos\theta
\right)
-
kx\cos^2\theta.
\end{equation}

On a

\begin{equation}
\frac{\partial\mathcal L}
{\partial\theta_x}
=
2x\theta_x
+
dx\sin\phi.
\end{equation}

Au bord $x=1$, la variation $\delta\theta(1)$ est libre. La condition
naturelle est donc

\begin{equation}
\left.
\frac{\partial\mathcal L}
{\partial\theta_x}
\right|_{x=1}
=0.
\end{equation}

On obtient

\begin{equation}
2\theta_x(1)
+
d\sin\phi
=0.
\end{equation}

Ainsi,

\begin{equation}
\boxed{
\theta_x(1)
=
-\frac{d}{2}\sin\phi.
}
\end{equation}

En revenant aux variables dimensionnelles,

\begin{equation}
\theta_x
=
R\theta_r,
\end{equation}

d'où

\begin{equation}
\boxed{
2A\theta_r(R)
+
D\sin\phi
=
0.
}
\end{equation}

Il s'agit de la condition de Brown généralisée dans le modèle considéré.

Dans le calcul numérique, cette condition n'est pas imposée comme une
contrainte. Elle doit être satisfaite par la configuration minimisante.

% ============================================================
\section{Équation d'Euler--Lagrange}
% ============================================================

Pour référence, la condition stationnaire dans le volume est

\begin{equation}
\frac{\partial\mathcal L}{\partial\theta}
-
\frac{\dd}{\dd x}
\left(
\frac{\partial\mathcal L}{\partial\theta_x}
\right)
=0.
\end{equation}

On obtient

\begin{equation}
\boxed{
\theta_{xx}
+
\frac{1}{x}\theta_x
-
\frac{\sin\theta\cos\theta}{x^2}
+
\frac{d}{x}\sin\phi\,\sin^2\theta
+
k\sin\theta\cos\theta
=
0.
}
\end{equation}

Cette équation est une condition de stationnarité de l'énergie.

Elle possède en général plusieurs solutions. Certaines peuvent être des
maxima ou des minima locaux de l'énergie.

C'est précisément pour cette raison que l'on préfère ici une minimisation
directe de la fonctionnelle pour rechercher l'état fondamental.

% ============================================================
\section{Discrétisation numérique}
% ============================================================

On introduit un maillage uniforme

\begin{equation}
x_i=ih,
\qquad
i=0,\ldots,N,
\end{equation}

avec

\begin{equation}
\boxed{
h=\frac{1}{N}.
}
\end{equation}

On note

\begin{equation}
\theta_i=\theta(x_i).
\end{equation}

La condition au centre est imposée par

\begin{equation}
\boxed{
\theta_0=0.
}
\end{equation}

Les variables optimisées sont donc

\begin{equation}
\boxed{
q=
(\theta_1,\theta_2,\ldots,\theta_N,\phi).
}
\end{equation}

Il y a ainsi $N+1$ variables d'optimisation.

La valeur $\theta_0$ n'est pas optimisée.

% ============================================================
\section{Approximation par éléments linéaires}
% ============================================================

Sur l'élément $[x_i,x_{i+1}]$, on approxime $\theta(x)$ par une fonction
linéaire :

\begin{equation}
\theta(x)
=
N_1(x)\theta_i
+
N_2(x)\theta_{i+1},
\end{equation}

avec les fonctions de forme

\begin{equation}
N_1(x)=\frac{x_{i+1}-x}{h},
\end{equation}

\begin{equation}
N_2(x)=\frac{x-x_i}{h}.
\end{equation}

La dérivée est constante dans l'élément :

\begin{equation}
\boxed{
\theta_x
=
\frac{\theta_{i+1}-\theta_i}{h}.
}
\end{equation}

% ============================================================
\section{Intégration de Gauss}
% ============================================================

L'intégrale sur chaque élément est calculée avec une quadrature de Gauss
à deux points.

Dans la coordonnée locale $\xi\in[-1,1]$, les deux points sont

\begin{equation}
\xi_1=-\frac{1}{\sqrt{3}},
\qquad
\xi_2=+\frac{1}{\sqrt{3}},
\end{equation}

avec les poids

\begin{equation}
w_1=w_2=1.
\end{equation}

La transformation est

\begin{equation}
x
=
\frac{x_i+x_{i+1}}{2}
+
\frac{h}{2}\xi.
\end{equation}

Le jacobien est

\begin{equation}
\dd x=\frac{h}{2}\dd\xi.
\end{equation}

La contribution d'un élément à l'énergie est donc évaluée comme

\begin{equation}
\mathcal E_i
\simeq
\frac{h}{2}
\sum_{a=1}^{2}
\mathcal L(x_{i,a},\theta_{i,a},\theta_x)
\end{equation}

et

\begin{equation}
\boxed{
\mathcal E
\simeq
\sum_{i=0}^{N-1}\mathcal E_i.
}
\end{equation}

% ============================================================
\section{Gradient analytique}
% ============================================================

L'élément essentiel de l'implémentation numérique est de fournir
explicitement le gradient de l'énergie à l'algorithme de minimisation.

On écrit

\begin{equation}
\nabla_q\mathcal E
=
\left(
\frac{\partial\mathcal E}{\partial\theta_1},
\ldots,
\frac{\partial\mathcal E}{\partial\theta_N},
\frac{\partial\mathcal E}{\partial\phi}
\right).
\end{equation}

% ------------------------------------------------------------
\subsection{Dérivée par rapport à $\theta$}
% ------------------------------------------------------------

On a

\begin{equation}
\frac{\partial\mathcal L}{\partial\theta}
=
\frac{2\sin\theta\cos\theta}{x}
+
d\sin\phi\cos(2\theta)
+
2kx\sin\theta\cos\theta.
\end{equation}

D'autre part,

\begin{equation}
\boxed{
\frac{\partial\mathcal L}
{\partial\theta_x}
=
2x\theta_x
+
dx\sin\phi.
}
\end{equation}

Puisque

\begin{equation}
\theta
=
N_1\theta_i+N_2\theta_{i+1},
\end{equation}

on a

\begin{equation}
\frac{\partial\theta}{\partial\theta_i}
=
N_1,
\qquad
\frac{\partial\theta}{\partial\theta_{i+1}}
=
N_2.
\end{equation}

Et

\begin{equation}
\frac{\partial\theta_x}{\partial\theta_i}
=
-\frac{1}{h},
\end{equation}

\begin{equation}
\frac{\partial\theta_x}{\partial\theta_{i+1}}
=
+\frac{1}{h}.
\end{equation}

Ainsi, pour un point de Gauss donné, la contribution au gradient du noeud
gauche est

\begin{equation}
\boxed{
\frac{\partial\mathcal L}{\partial\theta_i}
=
\frac{\partial\mathcal L}{\partial\theta}N_1
-
\frac{1}{h}
\frac{\partial\mathcal L}{\partial\theta_x}.
}
\end{equation}

Pour le noeud droit,

\begin{equation}
\boxed{
\frac{\partial\mathcal L}{\partial\theta_{i+1}}
=
\frac{\partial\mathcal L}{\partial\theta}N_2
+
\frac{1}{h}
\frac{\partial\mathcal L}{\partial\theta_x}.
}
\end{equation}

Les contributions provenant des éléments adjacents sont ensuite sommées.

% ------------------------------------------------------------
\subsection{Dérivée par rapport à $\phi$}
% ------------------------------------------------------------

La dépendance en $\phi$ est uniquement contenue dans le terme DMI :

\begin{equation}
d\sin\phi
\left(
x\theta_x+\sin\theta\cos\theta
\right).
\end{equation}

Par conséquent,

\begin{equation}
\boxed{
\frac{\partial\mathcal L}{\partial\phi}
=
d\cos\phi
\left(
x\theta_x
+
\sin\theta\cos\theta
\right).
}
\end{equation}

Le gradient global par rapport à $\phi$ est donc

\begin{equation}
\boxed{
\frac{\partial\mathcal E}{\partial\phi}
=
d\cos\phi
\int_0^1
\left(
x\theta_x
+
\sin\theta\cos\theta
\right)
\dd x.
}
\end{equation}

À l'état stationnaire, on doit avoir

\begin{equation}
\frac{\partial\mathcal E}{\partial\phi}=0.
\end{equation}

Pour une texture non triviale, cela conduit généralement à

\begin{equation}
\cos\phi=0,
\end{equation}

soit

\begin{equation}
\boxed{
\phi=\pm\frac{\pi}{2}.
}
\end{equation}

Le signe sélectionné dépend du signe de $D$ et de la texture obtenue.

% ============================================================
\section{Pourquoi utiliser le gradient analytique ?}
% ============================================================

Une minimisation numérique nécessite de calculer

\begin{equation}
\nabla_q\mathcal E.
\end{equation}

Une méthode naïve pourrait estimer chaque dérivée par différences finies :

\begin{equation}
\frac{\partial\mathcal E}{\partial q_j}
\simeq
\frac{
\mathcal E(q_j+\epsilon)
-
\mathcal E(q_j-\epsilon)
}
{2\epsilon}.
\end{equation}

Pour $N+1$ variables, cela nécessite de nombreuses évaluations de
l'énergie à chaque itération.

Lorsque $N$ augmente, le coût devient rapidement important.

Dans l'implémentation utilisée ici, le gradient est calculé analytiquement.
Une seule évaluation de l'énergie et du gradient est alors nécessaire à
chaque itération.

De plus, les opérations sur les éléments sont vectorisées avec
\texttt{NumPy}.

% ============================================================
\section{Algorithme de minimisation}
% ============================================================

On utilise l'algorithme

\begin{equation}
\boxed{
\texttt{L-BFGS-B}
}
\end{equation}

qui est une méthode quasi-newtonienne adaptée aux problèmes de grande
dimension et permettant d'imposer des bornes sur les variables.

Dans le programme, on impose

\begin{equation}
0\leq\theta_i\leq\pi,
\end{equation}

et

\begin{equation}
-\pi\leq\phi\leq\pi.
\end{equation}

La condition

\begin{equation}
\theta_0=0
\end{equation}

est imposée exactement en ne faisant pas de $\theta_0$ une variable
d'optimisation.

En revanche, aucune contrainte n'est imposée sur $\theta_N=\theta(R)$.

La condition de Brown doit donc émerger naturellement.

% ============================================================
\section{Recherche du minimum global}
% ============================================================

Une méthode de minimisation locale ne garantit pas à elle seule le
minimum global.

Il est donc utile de lancer l'algorithme depuis plusieurs configurations
initiales.

Par exemple,

\begin{equation}
\theta_{\mathrm{init}}(x)
=
\theta_{\max}
\tanh\left(\frac{x}{\delta}\right),
\end{equation}

avec différentes valeurs de $\delta$ et $\theta_{\max}$.

On peut également utiliser différentes valeurs initiales de $\phi$ :

\begin{equation}
\phi_{\mathrm{init}}
=
0,
\quad
+\frac{\pi}{2},
\quad
-\frac{\pi}{2},
\quad
\pi.
\end{equation}

On obtient alors plusieurs minima locaux

\begin{equation}
\mathcal E_1,
\mathcal E_2,\ldots
\end{equation}

et l'état retenu comme candidat à l'état fondamental est celui qui possède
l'énergie la plus faible :

\begin{equation}
\boxed{
\mathcal E_{\mathrm{GS}}
=
\min_j\mathcal E_j.
}
\end{equation}

% ============================================================
\section{Vérification numérique de la condition de Brown}
% ============================================================

Une fois la minimisation terminée, on calcule la dérivée au bord par une
différence finie :

\begin{equation}
\theta_x(1)
\simeq
\frac{\theta_N-\theta_{N-1}}{h}.
\end{equation}

On définit alors le résidu de Brown

\begin{equation}
\boxed{
B_{\mathrm{Brown}}
=
2\theta_x(1)
+
d\sin\phi.
}
\end{equation}

Une solution correctement minimisée doit vérifier

\begin{equation}
\boxed{
B_{\mathrm{Brown}}\simeq0.
}
\end{equation}

En unités dimensionnelles,

\begin{equation}
B_{\mathrm{Brown}}^{\mathrm{dim}}
=
2A\theta_r(R)
+
D\sin\phi.
\end{equation}

On doit donc également avoir

\begin{equation}
\boxed{
2A\theta_r(R)+D\sin\phi\simeq0.
}
\end{equation}

% ============================================================
\section{Contrôle de convergence en $N$}
% ============================================================

Pour vérifier que la solution numérique ne dépend pas du maillage, on peut
effectuer les calculs successivement pour

\begin{equation}
N=50,\quad100,\quad200,\quad400,\quad800.
\end{equation}

On doit vérifier la convergence de plusieurs grandeurs :

\begin{enumerate}
    \item l'énergie minimale ;
    \item la valeur de $\phi$ ;
    \item la valeur de $\theta(R)$ ;
    \item le profil $\theta(r)$ ;
    \item le résidu de Brown.
\end{enumerate}

On doit observer

\begin{equation}
E_N\longrightarrow E_\infty
\end{equation}

et

\begin{equation}
\theta_N(r)
\longrightarrow
\theta(r).
\end{equation}

Une convergence correcte doit également conduire à

\begin{equation}
B_{\mathrm{Brown}}\longrightarrow0.
\end{equation}

% ============================================================
\section{Composantes de l'aimantation}
% ============================================================

Une fois $\theta(r)$ et $\phi$ déterminés, les composantes de l'aimantation
sont

\begin{equation}
\boxed{
m_r(r)
=
\sin\theta(r)\cos\phi,
}
\end{equation}

\begin{equation}
\boxed{
m_\varphi(r)
=
\sin\theta(r)\sin\phi,
}
\end{equation}

et

\begin{equation}
\boxed{
m_z(r)
=
\cos\theta(r).
}
\end{equation}

Le programme représente notamment

\begin{equation}
\boxed{
m_z(r)=\cos\theta(r).
}
\end{equation}

Cette quantité permet de visualiser directement la composante de
l'aimantation parallèle à l'axe du cylindre.

% ============================================================
\section{Résumé de la méthode numérique}
% ============================================================

La procédure complète est donc

\begin{equation}
\boxed{
\begin{array}{c}
\text{Choix de }A,D,K_u,R,t
\\[1mm]
\downarrow
\\[1mm]
d=DR/A,\qquad k=K_uR^2/A
\\[1mm]
\downarrow
\\[1mm]
x=r/R
\\[1mm]
\downarrow
\\[1mm]
\theta_0=0
\\[1mm]
\downarrow
\\[1mm]
q=(\theta_1,\ldots,\theta_N,\phi)
\\[1mm]
\downarrow
\\[1mm]
\text{calcul de }\mathcal E(q)
\\[1mm]
\downarrow
\\[1mm]
\text{calcul analytique de }\nabla_q\mathcal E
\\[1mm]
\downarrow
\\[1mm]
\text{minimisation L-BFGS-B}
\\[1mm]
\downarrow
\\[1mm]
\text{plusieurs configurations initiales}
\\[1mm]
\downarrow
\\[1mm]
\text{sélection de l'énergie minimale}
\\[1mm]
\downarrow
\\[1mm]
\text{vérification de Brown}
\\[1mm]
\downarrow
\\[1mm]
\theta(r),\quad
m_z(r)=\cos\theta(r),\quad
\phi.
\end{array}
}
\end{equation}

L'intérêt principal de cette approche est qu'elle cherche directement un
minimum de la fonctionnelle d'énergie. Elle n'est donc pas limitée à la
recherche d'une solution particulière de l'équation d'Euler--Lagrange.

\end{document}
